\documentclass[../../../../main.tex]{subfiles}

\begin{document}

\section*{第1問}

\begin{proof}[解]
$f(x)=\dfrac{3}{4}x^2-3x+4=\dfrac{3}{4}(x-2)^2+1$ とおく。
$y=x$ と $y=f(x)$ の交点は $x=\dfrac{4}{3},4$ でグラフは右図。

\begin{center}
\begin{tikzpicture}[scale=1.1, >=stealth]
  \draw[->] (-0.3,0) -- (5,0) node[right] {$x$};
  \draw[->] (0,-0.3) -- (0,5) node[above] {$y$};
  \node[below left] at (0,0) {$O$};
  \draw[dashed] (0,0) -- (4.5,4.5) node[above right] {$y=x$};
  \draw[domain=0:4.6, smooth, thick] plot (\x, {0.75*(\x-2)*(\x-2)+1});
  \node[below] at (1.33,0) {$\frac{4}{3}$};
  \node[below] at (4,0) {$4$};
  \fill (1.33,1.33) circle (1.2pt);
  \fill (4,4) circle (1.2pt);
\end{tikzpicture}
\end{center}

したがって，$4<b$ の時，$f(b)>b$ で不適。次に $[a,b]$ に $2$ を含むかで場合分けする。

\textbf{$1^\circ\ a\le2\le b$ の時}

$[a,b]$ での $\min f=f(2)=1$ だから，$a=1$ である。$f(1)=\dfrac{7}{4}$ だから $b=4$ のみ適する。

\textbf{$2^\circ\ a\ge2$ の時}

$[a,b]$ では $f(a)\le f(x)\le f(b)$ だから，$f(a)=a$ かつ $f(b)=b$ とならば良いが，図からこれをみたす $(a,b)$ はない。

\textbf{$3^\circ\ b<2$ の時}

$[a,b]$ では $f(b)\le f(x)\le f(a)$ だから $f(b)=a$ かつ $f(a)=b$ となる。
\[
f(b)-f(a)=\frac{3}{4}(b^2-a^2)-3(b-a)=(a-b)
\]
$a\ne b$ から
\[
\frac{3}{4}(a+b)-3=-1 \quad\therefore\ a+b=\frac{8}{3}
\]
これを $f(a)=b$ に代入すると $a=b=\dfrac{4}{3}$ となり，$a=b$ に反し矛盾。

以上から，もとめるのは $(a,b)=(1,4)$ である。
\end{proof}

\newpage

\section*{第2問}

\begin{proof}[解]
$|\alpha|=k$ とおく。

\begin{enumerate}
  \item[(1)] $k\ne0,1$ とする。題意から $x=\alpha,\alpha^2,\alpha^4,\cdots,\alpha^{2^n},\cdots$ が $f(x)=0$ の解であり，$k\ne0,1$ から，これらは全て異なる。これは $f(x)$ が3次であることに反し矛盾。以上から，$k=0,1$ である。

  \item[(2)] $a,b,c\in\mathbb{R}$ として，$f(x)=x^3+ax^2+bx+c$ とおく。$f(x)=0$ の3解を $\alpha,\beta,\gamma$ とおく。

  \textbf{$1^\circ$ $\alpha,\beta,\gamma\in\mathbb{R}$ の時}

  $\alpha\ne\beta,\beta\ne\gamma,\gamma\ne\alpha$ 及び (1) から，$\{\alpha,\beta,\gamma\}=\{0,\pm1\}$ となり，この時，十分である。
  \[
  f(x)=x^3-x \quad \cdots \text{①}
  \]

  \textbf{$2^\circ$ $\alpha,\beta,\gamma$ のうち1つのみ実数の時}

  対称性から，$\alpha\in\mathbb{R}$ として良い。この時，$\gamma=\bar\beta$ となる（$\because a,b,c\in\mathbb{R}$）。又題意から，
  \[
  \alpha=0,1
  \]
  である。（(1) 及び，$\alpha=-1$ ならば $\alpha^2=1$ が解になるので）

  \textbf{(i) $\alpha=0$}

  題意から，$\beta^2$ も $f(x)=0$ の解だから，
  \[
  \beta^2=0,\beta,\bar\beta
  \]
  $\beta\ne0,1$ より，
  \[
  \beta^2=\bar\beta \quad \cdots (\star)
  \]
  ここで，$e(\theta)=\cos\theta+i\sin\theta$ として，$\beta=e(\theta)$ とすると，$(0\le\theta<2\pi)$，$(\star)$ から，（$\because(1)$）
  \[
  e(2\theta)=e(-\theta) \quad\therefore\ 2\theta=-\theta+2n\pi\ (n\in\mathbb{Z}) \quad\therefore\ \theta=\frac{2}{3}n\pi
  \]
  $0\le\theta<2\pi$ から，$n=0,1,2$ として，
  \[
  \theta=0,\frac{2\pi}{3},\frac{4\pi}{3}
  \]
  このうち，$\theta=0$ は $\beta=1$ となって不適で，その他の時は十分である。この時
  \[
  f(x)=x^3+x^2+x \quad \cdots \text{②}
  \]
  となる。

  \textbf{(ii) $\alpha=1$}

  (i) と同じく，
  \[
  \beta^2=1,\beta,\bar\beta
  \]
  $\beta\ne0,1$ から，
  \[
  \beta^2=1,\bar\beta
  \]
  $\beta^2=\bar\beta$ の時，(i) と同じく，$\{\beta,\gamma\}=\left\{e\left(\dfrac{2}{3}\pi\right),e\left(\dfrac{4}{3}\pi\right)\right\}$ で，
  \[
  f(x)=(x-1)(x^2+x+1)=x^3-1 \quad \cdots \text{③}
  \]
  $\beta^2=1$ の時，$\beta=\pm1$ で，$\beta\in\mathbb{R}$ となり矛盾。

  $a,b,c\in\mathbb{R}$ から，以上で全ての場合がつくせる。よって①②③より，
  \[
  f(x)=x^3-x,\ x^3+x^2+x,\ x^3-1
  \]
\end{enumerate}
\end{proof}

\newpage

\section*{第3問}

\begin{proof}[解]
\begin{enumerate}
  \item[(1)] $w\in\mathbb{R} \iff w=\bar w$ だから $w=az^2+bz$ を代入して
  \[
  az^2+bz=a\bar z^2+b\bar z
  \iff a(z^2-\bar z^2)+b(z-\bar z)=0
  \iff \{a(z+\bar z)+b\}(z-\bar z)=0
  \]
  \[
  \iff z+\bar z=-\frac{b}{a}\ \text{or}\ z\in\mathbb{R}
  \iff x=-\frac{b}{2a}\ \text{or}\ y=0
  \]

  \item[(2)] $y=0$ の時，$w$ は $f(x)=ax^2+bx$ で表される。 $\cdots$①

  次に $x=-\dfrac{b}{2a}$ の時，
  \[
  az^2+bz = a\left(-\frac{b}{2a}+iy\right)^2+b\left(-\frac{b}{2a}+iy\right)
  = a\left(\frac{b^2}{4a^2}-\frac{by}{a}i-y^2\right)+byi-\frac{b^2}{2a}
  \]
  \[
  = -ay^2-\frac{b^2}{4a} \equiv g(y) \quad \cdots \text{②}
  \]
  である。

  \textbf{$1^\circ\ a>0$ の時}

  $f(x)\ge-\dfrac{b^2}{4a}$，$g(y)\le-\dfrac{b^2}{4a}$ でこの間の値はくまなくとるから，$w$ は任意の実数値をとる。

  \textbf{$2^\circ\ a<0$ の時}

  $f(x)\le-\dfrac{b^2}{4a}$，$g(y)\ge-\dfrac{b^2}{4a}$ で $1^\circ$ と同様に $w$ は任意の実数値をとる。

  以上から示された。
\end{enumerate}
\end{proof}

\newpage

\section*{第4問}

\begin{proof}[解]
対称性から，立体のうち $y\ge0$ の部分の体積 $V_1$，もとめる体積 $V_2$ として，
\[
V_2=2V_1 \quad \cdots \text{①}
\]
である。接線は $y$ 軸平行でないので，$m$ を傾きとして
\[
\ell: y=m(x-2a)
\]
とおける。これが楕円と接するので，$X=x/a,\ Y=y/b$ とおき直した座標で，
\[
bY=m(aX-2a)
\]
と $X^2+Y^2=1$ が接する。したがって
\[
\frac{|2am|}{\sqrt{(am)^2+b^2}}=1
\]
各辺0以上から2乗して，
\[
4a^2m^2=a^2m^2+b^2 \quad\therefore\ m=\pm\frac{\sqrt3}{3}\frac{b}{a}
\]
図の接線の傾きは負だから，$a,b>0$ より複号負をとって，
\[
\ell: y=-\frac{\sqrt3}{3}\frac{b}{a}(x-2a)
\]
である。この時接点 $P\left(\dfrac{a}{2},\dfrac{\sqrt3}{2}b\right)$ となる。したがって，
\[
V_1 = (\text{円錐部分}) - (\text{楕円部分}) \quad \cdots \text{②}
\]
において，
\[
(\text{円錐部分}) = \frac{1}{3}\pi(2a)^2\frac{2\sqrt3}{3}b-\frac{1}{3}\pi\left(\frac{1}{2}a\right)^2\left(\frac{2\sqrt3}{3}b-\frac{\sqrt3}{2}b\right)
= \frac{1}{3}\pi\left[\frac{8\sqrt3}{3}a^2b-\frac{\sqrt3}{24}a^2b\right] = \frac{7\sqrt3}{8}\pi a^2b
\]
\[
(\text{楕円部分}) = \pi\int_0^{\frac{\sqrt3}{2}b}a^2\left(1-\frac{y^2}{b^2}\right)dy=a^2\pi\left[y-\frac{1}{3b^2}y^3\right]_0^{\frac{\sqrt3}{2}b}=\frac{3\sqrt3}{8}\pi a^2b
\]
を②に代入して
\[
V_1=\frac{7-3}{8}\sqrt3\pi a^2b=\frac{\sqrt3}{2}\pi a^2b
\]
①から
\[
V_2=\sqrt3\pi a^2b
\]

\begin{center}
\begin{tikzpicture}[scale=0.9, >=stealth]
  \draw[->] (-0.3,0) -- (4.5,0) node[right] {$x$};
  \draw[->] (0,-0.3) -- (0,3) node[above] {$y$};
  \draw[thick] (0,0) ellipse (1.5cm and 1.3cm);
  \node[below] at (1.5,0) {$a$};
  \node[left] at (0,1.3) {$b$};
  \fill (3,0) circle (1pt) node[below] {$2a$};
  \draw[thick] (0,1.13) -- (3,0);
  \fill (0.75,1.13) circle (1pt) node[above] {$P$};
  \begin{scope}
    \clip (0,0) rectangle (3,1.3);
    \fill[pattern=north east lines, pattern color=gray!50] (1.5,0) -- (3,0) -- (0.75,1.13) arc (37:0:1.5 and 1.3) -- cycle;
  \end{scope}
\end{tikzpicture}
\end{center}
\end{proof}

\newpage

\section*{第5問}

\begin{proof}[解]
$f(x)=e^x-ax$ とおく。$f'(x)=e^x-a$ から，$a$ によって増減は以下のようになる。

\textbf{$1^\circ\ a\le1$ の時}

$f'(x)\ge0$ から，$f(x)$ は単調増加で，$f(0)=1,f(1)=e-a$ だから，
\[
\max|f(x)|=2 \iff e-a=2 \iff a=e-2
\]
でこれは $a\le1$ をみたす。

\textbf{$2^\circ\ 1\le a\le e$ の時}

$f'(x)=0$ なる $x$ が $0\le x\le1$ にただ1つ存在し，下表のようになる。（これを $\alpha$ とおく）

\begin{center}
\begin{tabular}{c|c|c|c}
$x$ & $0$ & $\alpha$ & $1$ \\
\hline
$f'$ & & $-$\ \ $0$\ \ $+$ & \\
\hline
$f$ & $1$ & $\searrow$\ \ \ $\nearrow$ & $e-a$
\end{tabular}
\end{center}
（$f'(\alpha)=0$ から，$e^\alpha=a$ で $1\le a\le e$ から $0\le\alpha\le1$）

この時，$\max|f(x)|=2$ の時，$|e-a|<2$ から $f(\alpha)=-1$ であるが，
\[
f(\alpha)=e^\alpha-a\alpha=e^\alpha(1-\alpha)>0
\]
から矛盾。

\textbf{$3^\circ\ e\le a$ の時}

$f'(x)\le0$ から $f(x)$ は単調減少で，$f(0)=1,f(1)=e-a$ だから，$f(1)\le0$ とあわせて，
\[
\max|f(x)|=2\iff e-a=-2\iff a=e+2
\]
でこれは $a\ge e$ をみたす。

以上から，もとめるのは $a=e\pm2$
\end{proof}

\newpage

\section*{第6問}

\begin{proof}[解]
$F(a)=\displaystyle\int_0^{\pi/2}|\sin x-a\cos x|dx$，$S=\sin x,\ C=\cos x,\ f(x)=S-aC$ とおく。$f'(x)=C+aS$

$a\le0$ の時，$f(x)\ge0$ とでき，
\[
F(a)=\int_0^{\pi/2}S-a\int_0^{\pi/2}C\,dx=1-a
\]
これは $a$ の単調減少関数であるから，連続性も考えて，$0\le a$ で考えれば良い。この時
$f(0)=-a\le0$，$f(\pi/2)=1$，$f'(x)\ge0$ から，$[0,\pi/2]$ に $f(\alpha)=0$ なる $\alpha$ がただ1つあって，
\[
F(a)=-\int_0^\alpha f(x)dx+\int_\alpha^{\pi/2}f(x)dx
= -\int_0^\alpha S\,dx+\int_0^\alpha aC\,dx+\int_\alpha^{\pi/2}S\,dx-\int_\alpha^{\pi/2}aC\,dx \quad \cdots \text{①}
\]
とかける。
\[
F'(a) = -\sin\alpha\cdot\alpha'+a\cos\alpha\cdot\alpha'+\int_0^\alpha C\,dx-\sin\alpha\cdot\alpha'+a\cos\alpha\cdot\alpha'-\int_\alpha^{\pi/2}C\,dx
\]
\[
=2(a\cos\alpha-\sin\alpha)\alpha'+2\sin\alpha-1 = 2\sin\alpha-1 \quad (\because f(\alpha)=0)
\]
から下表をうる。

\begin{center}
\begin{tabular}{c|c|c|c|c}
$a$ & $0$ & & & $+\infty$ \\
\hline
$\alpha$ & $0$ & \multicolumn{2}{c|}{$\nearrow$} & $(\pi/2)$ \\
\hline
$F'$ & & $-$ & $0$ & $+$ \\
\hline
$F$ & & $\searrow$ & $\dfrac{\pi}{6}$ & $\nearrow$
\end{tabular}
\end{center}

したがって，$\alpha=\dfrac{\pi}{6}$ で $F(a)$ は最小で，この時，$f(\alpha)=0$ より，$a=\dfrac{\sqrt3}{3}$ でこの時，①から
\[
F\left(\frac{\sqrt3}{3}\right) = -\int_0^{\pi/6}\left(S-\frac{\sqrt3}{3}C\right)dx+\int_{\pi/6}^{\pi/2}\left(S-\frac{\sqrt3}{3}C\right)dx
\]
\[
= -\left[-C-\frac{\sqrt3}{3}S\right]_0^{\pi/6}+\left[-C-\frac{\sqrt3}{3}S\right]_{\pi/6}^{\pi/2}
= \left(\frac{\sqrt3}{2}+\frac{\sqrt3}{6}-1\right)-\left(\frac{\sqrt3}{3}-\frac{\sqrt3}{2}-\frac{\sqrt3}{6}\right)
\]
\[
= \sqrt3-1
\]
となる。
\end{proof}

\end{document}
